\subsubsection{Compact journey choices and transfer-event accounting}
\label{sec:reduced-od-journey-choices}

Phase~5 converts the timetable labels into complete, bounded passenger-journey
alternatives in
\path{public_transportation.preprocessing.reduced_od.journey_choices}.  This
step remains preprocessing: it neither changes demand nor estimates route
shares.  One choice cell is identified by physical origin, physical
destination, and the half-open time period containing the journey's first
boarding.  A public-transport alternative must contain at least one scheduled
vehicle leg; walking-only reachability cannot silently become a transit-demand
cell.

For a journey with $L$ vehicle legs, the event sequence contains exactly $2L$
records.  The first vehicle entry is labeled
\texttt{first\_boarding}, the final exit is labeled
\texttt{final\_alighting}, and each of the $L-1$ internal transfers contributes
the consecutive pair
\[
\bigl(\texttt{transfer\_alighting},
      \texttt{transfer\_boarding}\bigr).
\]
The two transfer events may use different physical stops when an explicit
footpath connects the arrival and departure locations, but the subsequent
boarding time cannot precede the preceding alighting.  Thus every journey has
the same number of boarding and alighting events, and transfer boardings remain
distinguishable from exogenous journey origins.  Every event, not merely the
journey origin, is assigned to exactly one declared half-open period; this
allows one journey to cross time-period boundaries without relabeling its
origin.

Candidate labels are deduplicated by a SHA-256 identity containing the query,
OD pair, first-boarding period, and complete scheduled leg sequence.  Within a
cell they are ranked deterministically by the generalized time
\[
c_j=T_{\mathrm{elapsed},j}
    +\gamma_K K_j
    +(w_{\mathrm{walk}}-1)T_{\mathrm{walk},j},
\]
followed by arrival time, transfers, and identifier.  The subtraction in the
walking term is important because elapsed time already contains walking: a
walking weight of one therefore leaves ordinary elapsed time unchanged.  Only
the configured maximum number of alternatives is retained, and diagnostics
report candidate, retained, and pruned counts as well as the largest
pre-pruning cell.
The result fingerprint additionally covers every derived feature, route-pattern
identifier, event kind, event period, event stop and time, and fixed share; a
change in downstream measurement semantics therefore changes the artifact
identity even when the scheduled leg sequence itself is unchanged.

Initial alternative shares are fixed metadata for the later estimator.  With
temperature $\tau$ and optional positive route-pattern initialization weight
$r_j$ for the first boarded pattern, they are
\[
s_j=
\frac{r_j\exp[-(c_j-c_{\min})/\tau]}
     {\sum_k r_k\exp[-(c_k-c_{\min})/\tau]}.
\]
Missing route-pattern weights default to one; supplied unknown, nonfinite, or
nonpositive weights are rejected.  This is the explicit interface through
which the route-level initialization may influence network choices.  It does
not reinterpret route-level IPF output as journey demand, and the shares remain
constant within a subsequent statistical fit.

The Phase-5 benchmark measures choice construction after RAPTOR has already
finished, so routing and choice-building time are not conflated.  Across the
25-, 50-, 100-, and 200-stop synthetic cases it processed approximately
70,000--77,000 choice cells per second.  The fixture retained one Pareto
alternative per cell and used approximately 607--655 bytes of full canonical
feature, leg, event, and share payload per retained alternative.  Multi-route unit tests separately exercise
two-alternative cells, transfer-event pairing, route-weighted shares,
deterministic pruning, and journeys crossing period boundaries.  The benchmark
values are operational diagnostics, not performance thresholds.

